% GENERATED FILE. Edit canonical lessons, then run npm run generate:books.
% canonical-source-hash: fnv1a64:dc544d69c22e9a0d
% canonical-lessons: JA-C503-utagoe, JA-C503-senritsu, JA-C503-onpu, JA-C503-ongakuka, JA-C503-mau, JA-R503-review-498-503

\chapter{Melodies and Dances}
\label{ch:ja-melodies-and-dances}
\hlchaptermodality{503}

\begin{chapteropening}
\textbf{By the end of this chapter:} \emph{I can talk about music and dance.}

Complete the last lesson of chapter 503: I can talk about music and dance.
\end{chapteropening}

\section[\texorpdfstring{\ja{うたごえ} (utagoe)}{うたごえ (utagoe)}]{\ja{うたごえ} (utagoe) — a singing voice; the sound of singing}

\label{lesson:JA-C503-utagoe}



\begin{quote}
Before the new one: say the Japanese for an actress, then the Japanese for to play (a role).
\end{quote}

\subsection*{You'll want to know: \ja{うたごえ}}
\textbf{\ja{うたごえ}} — \emph{utagoe} — \textquotedblleft{}a singing voice; the sound of singing\textquotedblright{}.

\begin{quote}
\textbf{\ja{となりの} \ja{へやから} \ja{うたごえが} \ja{しました。}} — \emph{tonari no heya kara utagoe ga shimashita} — \textquotedblleft{}There was singing from the next room.\textquotedblright{}
\end{quote}

\begin{cousinweb}
\textbf{\ja{うたごえ}} (\emph{utagoe}) is a native Japanese compound (wago), inherited rather than borrowed from Chinese: built from \emph{uta}, \textquotedblleft{}song\textquotedblright{} and \emph{goe}, \textquotedblleft{}voice\textquotedblright{} (JMdict, KANJIDIC2).
\end{cousinweb}

\begin{grammarlens}[title={What you've built}]
One more word for describing what you lived and felt.
\end{grammarlens}

\subsection*{Your turn}
\begin{itemize}
\raggedright
  \item \emph{Say it:} \emph{utagoe}
  \item \emph{Say it:} \emph{utagoe}, once more
  \item \emph{Say it:} \emph{tonari no heya kara utagoe ga shimashita} — There was singing from the next room.
  \item \emph{Recall:} say the Japanese for an actress, then the Japanese for to play (a role), then say \emph{utagoe} again
\end{itemize}

\begin{tcolorbox}[breakable,colback=teal!4,colframe=teal!35!black,title={Before you move on}]
What does \ja{うたごえ} mean? (\textquotedblleft{}A singing voice; the sound of singing\textquotedblright{}.) Say it once more.
\end{tcolorbox}
\section[\texorpdfstring{\ja{せんりつ} (senritsu)}{せんりつ (senritsu)}]{\ja{せんりつ} (senritsu) — a melody; a tune}

\label{lesson:JA-C503-senritsu}



\begin{quote}
Before the new one: say the Japanese for to play (a role), then the Japanese for a singing voice.
\end{quote}

\subsection*{You'll want to know: \ja{せんりつ}}
\textbf{\ja{せんりつ}} — \emph{senritsu} — \textquotedblleft{}a melody; a tune\textquotedblright{}.

\begin{quote}
\textbf{\ja{この} \ja{うたの} \ja{せんりつは} \ja{うつくしいです。}} — \emph{kono uta no senritsu wa utsukushii desu} — \textquotedblleft{}The melody of this song is beautiful.\textquotedblright{}
\end{quote}

\begin{cousinweb}
\textbf{\ja{せんりつ}} (\emph{senritsu}) is Sino-Japanese (kango): two Chinese characters read with their on'yomi, readings borrowed from Chinese. They are \emph{sen}, \textquotedblleft{}rotation\textquotedblright{} (Middle Chinese \emph{zjwen}) and \emph{ritsu}, \textquotedblleft{}rhythm\textquotedblright{} (\emph{lwit}) (KANJIDIC2; Middle Chinese in Baxter's transcription).
\end{cousinweb}

\begin{grammarlens}[title={What you've built}]
One more word for describing what you lived and felt.
\end{grammarlens}

\subsection*{Your turn}
\begin{itemize}
\raggedright
  \item \emph{Say it:} \emph{senritsu}
  \item \emph{Say it:} \emph{senritsu}, once more
  \item \emph{Say it:} \emph{kono uta no senritsu wa utsukushii desu} — The melody of this song is beautiful.
  \item \emph{Recall:} say the Japanese for to play (a role), then the Japanese for a singing voice, then say \emph{senritsu} again
\end{itemize}

\begin{tcolorbox}[breakable,colback=teal!4,colframe=teal!35!black,title={Before you move on}]
What does \ja{せんりつ} mean? (\textquotedblleft{}A melody; a tune\textquotedblright{}.) Say it once more.
\end{tcolorbox}
\section[\texorpdfstring{\ja{おんぷ} (onpu)}{おんぷ (onpu)}]{\ja{おんぷ} (onpu) — a musical note}

\label{lesson:JA-C503-onpu}



\begin{quote}
Before the new one: say the Japanese for a singing voice, then the Japanese for a melody.
\end{quote}

\subsection*{You'll want to know: \ja{おんぷ}}
\textbf{\ja{おんぷ}} — \emph{onpu} — \textquotedblleft{}a musical note\textquotedblright{}.

\begin{quote}
\textbf{\ja{この} \ja{うたは} \ja{おんぷが} \ja{おおいです。}} — \emph{kono uta wa onpu ga ōi desu} — \textquotedblleft{}This song has a lot of notes.\textquotedblright{}
\end{quote}

\begin{cousinweb}
\textbf{\ja{おんぷ}} (\emph{onpu}) is Sino-Japanese (kango): two Chinese characters read with their on'yomi, readings borrowed from Chinese. They are \emph{on}, \textquotedblleft{}sound\textquotedblright{} (Middle Chinese \emph{'im}) and \emph{pu}, \textquotedblleft{}token\textquotedblright{} (\emph{bju}) (KANJIDIC2; Middle Chinese in Baxter's transcription).
\end{cousinweb}

\begin{grammarlens}[title={What you've built}]
One more word for describing what you lived and felt.
\end{grammarlens}

\subsection*{Your turn}
\begin{itemize}
\raggedright
  \item \emph{Say it:} \emph{onpu}
  \item \emph{Say it:} \emph{onpu}, once more
  \item \emph{Say it:} \emph{kono uta wa onpu ga ōi desu} — This song has a lot of notes.
  \item \emph{Recall:} say the Japanese for a singing voice, then the Japanese for a melody, then say \emph{onpu} again
\end{itemize}

\begin{tcolorbox}[breakable,colback=teal!4,colframe=teal!35!black,title={Before you move on}]
What does \ja{おんぷ} mean? (\textquotedblleft{}A musical note\textquotedblright{}.) Say it once more.
\end{tcolorbox}
\section[\texorpdfstring{\ja{おんがくか} (ongakuka)}{おんがくか (ongakuka)}]{\ja{おんがくか} (ongakuka) — a musician}

\label{lesson:JA-C503-ongakuka}



\begin{quote}
Before the new one: say the Japanese for a melody, then the Japanese for a musical note.
\end{quote}

\subsection*{You'll want to know: \ja{おんがくか}}
\textbf{\ja{おんがくか}} — \emph{ongakuka} — \textquotedblleft{}a musician\textquotedblright{}.

\begin{quote}
\textbf{\ja{おじさんは} \ja{おんがくかです。}} — \emph{ojisan wa ongakuka desu} — \textquotedblleft{}My uncle is a musician.\textquotedblright{}
\end{quote}

\begin{cousinweb}
\textbf{\ja{おんがくか}} (\emph{ongakuka}) is Sino-Japanese (kango): three Chinese characters read with their on'yomi, readings borrowed from Chinese. They are \emph{on}, \textquotedblleft{}sound\textquotedblright{} (Middle Chinese \emph{'im}), \emph{gaku}, \textquotedblleft{}music\textquotedblright{} (\emph{ngaewk}) and \emph{ka}, \textquotedblleft{}expert\textquotedblright{} (\emph{kae}) (KANJIDIC2; Middle Chinese in Baxter's transcription).
\end{cousinweb}

\begin{grammarlens}[title={What you've built}]
One more word for describing what you lived and felt.
\end{grammarlens}

\subsection*{Your turn}
\begin{itemize}
\raggedright
  \item \emph{Say it:} \emph{ongakuka}
  \item \emph{Say it:} \emph{ongakuka}, once more
  \item \emph{Say it:} \emph{ojisan wa ongakuka desu} — My uncle is a musician.
  \item \emph{Recall:} say the Japanese for a melody, then the Japanese for a musical note, then say \emph{ongakuka} again
\end{itemize}

\begin{tcolorbox}[breakable,colback=teal!4,colframe=teal!35!black,title={Before you move on}]
What does \ja{おんがくか} mean? (\textquotedblleft{}A musician\textquotedblright{}.) Say it once more.
\end{tcolorbox}
\section[\texorpdfstring{\ja{まう} (mau)}{まう (mau)}]{\ja{まう} (mau) — to flutter; to dance (in the air)}

\label{lesson:JA-C503-mau}



\begin{quote}
Before the new one: say the Japanese for a musical note, then the Japanese for a musician.
\end{quote}

\subsection*{You'll want to know: \ja{まう}}
\textbf{\ja{まう}} — \emph{mau} — \textquotedblleft{}to flutter; to dance (in the air)\textquotedblright{}.

\begin{quote}
\textbf{\ja{ゆきが} \ja{かぜに} \ja{まっています。}} — \emph{yuki ga kaze ni matteimasu} — \textquotedblleft{}The snow is dancing in the wind.\textquotedblright{}
\end{quote}

Polite \textbf{\ja{まいます}}, \emph{maimasu}; a godan verb. Past (polite): \textbf{\ja{ゆきが} \ja{まいました。}}, \emph{yuki ga maimashita}.

\begin{cousinweb}
\textbf{\ja{まう}} (\emph{mau}) is a native Japanese word (wago), inherited rather than borrowed from Chinese; its kanji was attached later as a kun'yomi, and its older origin is not settled (JMdict, KANJIDIC2).
\end{cousinweb}

\begin{grammarlens}[title={What you've built}]
One more word for describing what you lived and felt.
\end{grammarlens}

\subsection*{Your turn}
\begin{itemize}
\raggedright
  \item \emph{Say it:} \emph{mau}
  \item \emph{Say it:} \emph{mau}, once more
  \item \emph{Say it:} \emph{yuki ga kaze ni matteimasu} — The snow is dancing in the wind.
  \item \emph{Recall:} say the Japanese for a musical note, then the Japanese for a musician, then say \emph{mau} again
\end{itemize}

\begin{tcolorbox}[breakable,colback=teal!4,colframe=teal!35!black,title={Before you move on}]
What does \ja{まう} mean? (\textquotedblleft{}To flutter; to dance (in the air)\textquotedblright{}.) Say it once more.
\end{tcolorbox}
\section[(dialogue)]{Looking back: from A Day Off to Melodies and Dances}

\label{lesson:JA-R503-review-498-503}



\begin{quote}
Say the Japanese for a day off and the winter holidays.
\end{quote}

\subsection*{Your turn}
Say these aloud:

\begin{itemize}
\raggedright
  \item a day off, the winter holidays, a nap, nodding off, paid leave
  \item an amusement park, a zoo, an aquarium, a school outing, a pleasure trip
  \item a toy, skipping rope, a running race (between children), a riddle, a slide (in a playground)
  \item moon viewing, flower viewing, autumn leaves, a wind chime, to scatter (something)
  \item a theatre, drama, a stage actor, an actress, to play (a role)
  \item a singing voice, a melody, a musical note, a musician, to flutter
\end{itemize}

\begin{tcolorbox}[breakable,colback=teal!4,colframe=teal!35!black,title={Before you move on}]
A day off? (\textbf{\ja{きゅうじつ}}, \emph{kyūjitsu}.) The winter holidays? (\textbf{\ja{ふゆやすみ}}, \emph{fuyuyasumi}.) A nap? (\textbf{\ja{ひるね}}, \emph{hirune}.) Nodding off? (\textbf{\ja{いねむり}}, \emph{inemuri}.) Paid leave? (\textbf{\ja{ゆうきゅう}}, \emph{yūkyū}.) An amusement park? (\textbf{\ja{ゆうえんち}}, \emph{yūenchi}.) A zoo? (\textbf{\ja{どうぶつえん}}, \emph{dōbutsuen}.) An aquarium? (\textbf{\ja{すいぞくかん}}, \emph{suizokukan}.) A school outing? (\textbf{\ja{えんそく}}, \emph{ensoku}.) A pleasure trip? (\textbf{\ja{こうらく}}, \emph{kōraku}.) A toy? (\textbf{\ja{おもちゃ}}, \emph{omocha}.) Skipping rope? (\textbf{\ja{なわとび}}, \emph{nawatobi}.) A running race (between children)? (\textbf{\ja{かけっこ}}, \emph{kakekko}.) A riddle? (\textbf{\ja{なぞなぞ}}, \emph{nazonazo}.) A slide (in a playground)? (\textbf{\ja{すべりだい}}, \emph{suberidai}.) Moon viewing? (\textbf{\ja{つきみ}}, \emph{tsukimi}.) Flower viewing? (\textbf{\ja{はなみ}}, \emph{hanami}.) Autumn leaves? (\textbf{\ja{もみじ}}, \emph{momiji}.) A wind chime? (\textbf{\ja{ふうりん}}, \emph{fūrin}.) To scatter (something)? (\textbf{\ja{ちらす}}, \emph{chirasu}.) A theatre? (\textbf{\ja{げきじょう}}, \emph{gekijō}.) Drama? (\textbf{\ja{えんげき}}, \emph{engeki}.) A stage actor? (\textbf{\ja{やくしゃ}}, \emph{yakusha}.) An actress? (\textbf{\ja{じょゆう}}, \emph{joyū}.) To play (a role)? (\textbf{\ja{えんじる}}, \emph{enjiru}.) A singing voice? (\textbf{\ja{うたごえ}}, \emph{utagoe}.) A melody? (\textbf{\ja{せんりつ}}, \emph{senritsu}.) A musical note? (\textbf{\ja{おんぷ}}, \emph{onpu}.) A musician? (\textbf{\ja{おんがくか}}, \emph{ongakuka}.) To flutter? (\textbf{\ja{まう}}, \emph{mau}.)
\end{tcolorbox}
